\section*{Chapitre \ref{ch:b3:bioinorganic} --- Chimie bioinorganique}

\begin{solution}{exo:b3:bioinorganic:1}
$\log(\theta/(1 - \theta))$ passe de $\log 0.25 = -0.602$ à $\log 4 = +0.602$ tandis que $\log p$ augmente de
$\log(5.9/2.1) = 0.449$~: $n = 1.204/0.449 = 2.7$.
\end{solution}

\begin{solution}{exo:b3:bioinorganic:2}
$\ce{Mn^2+} < \ce{Fe^2+} < \ce{Co^2+} < \ce{Ni^2+} < \ce{Cu^2+} > \ce{Zn^2+}$. Le zinc ($d^{10}$) n'a ni stabilisation de champ des ligands ni
distorsion de Jahn--Teller~: il retombe au-dessous du cuivre.
\end{solution}

\begin{solution}{exo:b3:bioinorganic:3}
$\ce{Ca^2+}$ et $\ce{Fe^3+}$, durs~: carboxylates (et phénolates pour le fer). $\ce{Cu+}$ et $\ce{Hg^2+}$, mous~:
thiolates de cystéine. $\ce{Zn^2+}$, intermédiaire~: histidine et cystéine.
\end{solution}

\begin{solution}{exo:b3:bioinorganic:4}
$\ce{2H2O -> O2 + 4H+ + 4e-}$~: quatre électrons, quatre photons par $\ce{O2}$~; quatre moles de photons par
mole de $\ce{O2}$.
\end{solution}

\begin{solution}{exo:b3:bioinorganic:5}
Myoglobine~: $\theta = 13/13.37 = 0.972$ et $5/5.37 = 0.931$~: elle cède $(0.972 - 0.931)/0.972 = 4$\,\% de son
dioxygène. Hémoglobine~: $\theta = 0.972$ et $0.724$~: elle en cède 26\,\%.
\end{solution}

\begin{solution}{exo:b3:bioinorganic:6}
$[\mathrm{HbCO}]/[\mathrm{HbO_2}] = 230 \times \num{50e-6}/0.2095 = 0.055$~; fraction $0.055/1.055 = 5.2$\,\%.
\end{solution}

\begin{solution}{exo:b3:bioinorganic:7}
Sulfures $-2$ chacun, cystéinates $-1$ chacun. $\ce{[Fe2S2(SR)4]^2-}$~: somme des fers $-2 + 4 + 4 = +6$~: deux
Fe(III). $\ce{[Fe2S2(SR)4]^3-}$~: $+5$, un Fe(III) et un Fe(II) (valence mixte). $\ce{[Fe4S4(SR)4]^2-}$~:
$-2 + 8 + 4 = +10$, deux Fe(III) et deux Fe(II), soit $+2.5$ en moyenne par fer.
\end{solution}

\begin{solution}{exo:b3:bioinorganic:8}
$\ce{[PtCl2(NH3)2] + H2O -> [PtCl(H2O)(NH3)2]+ + Cl-}$. La forte concentration en chlorure du sang
ramène l'équilibre en arrière~; dans les cellules, le chlorure est bien
moins concentré, et l'aquacomplexe se forme. Le platine se lie à l'atome N7
de la guanine. Dans l'isomère \textit{trans}, les deux sites \omterm{def:b3:complex-mechanisms:labile}{labiles} sont
opposés et ne peuvent pas atteindre deux bases adjacentes d'un même brin.
\end{solution}

\begin{solution}{exo:b3:bioinorganic:9}
$1/T_1 = 1/\qty{1.2}{s} + \qty{4.0}{L.mmol^{-1}.s^{-1}} \times \qty{0.10}{mmol/L} = 0.833 + 0.400 = \qty{1.23}{s^{-1}}$~: $T_1 = \qty{0.81}{s}$.
\end{solution}

\begin{solution}{exo:b3:bioinorganic:10}
$K = K_{\mathrm{PbL}}/K_{\mathrm{CaL}} = 10^{18.0 - 10.7} = 10^{7.3} = \num{2e7}$~: le plomb déplace totalement le calcium. Le ligand
libre fixerait aussi le calcium du sang et abaisserait dangereusement sa
concentration~; le complexe calcique s'échange avec le plomb et laisse le
taux de calcium inchangé.
\end{solution}

\begin{solution}{exo:b3:bioinorganic:11}
Bien au-dessous de $K_M$, la réaction est du premier ordre~: $k' = (k_{\mathrm{cat}}/K_M)[\mathrm E]$, $t_{1/2} = \ln 2/k'$. Anhydrase
carbonique~: $k' = \qty{100}{s^{-1}}$, $t_{1/2} = \qty{6.9}{ms}$. \omterm{def:b3:complex-kinetics:enzyme}{Enzyme} typique~: $k' = \qty{0.10}{s^{-1}}$, $t_{1/2} = \qty{6.9}{s}$, mille
fois plus long, trop lent pour la seconde que passe un globule rouge dans
un capillaire pulmonaire.
\end{solution}

\begin{solution}{exo:b3:bioinorganic:12}
La moitié des sites occupés par CO signifie $[\mathrm{HbCO}] = [\mathrm{HbO_2}]$~: $p(\ce{CO}) = p(\ce{O2})/M = 0.2095/230 = \qty{9.1e-4}{bar}$, soit environ
\qty{910}{ppm}.
\end{solution}

\begin{solution}{pb:b3:bioinorganic:1}
\textbf{1.} $\theta = 13^{2.7}/(3.5^{2.7} + 13^{2.7}) = 0.972$.

\textbf{2.} $\theta(\qty{5.0}{kPa}) = 0.724$.

\textbf{3.} $5.0/(0.37 + 5.0) = 0.931$.

\textbf{4.} Son $p_{50}$ est dix fois plus faible~: aux pressions d'un muscle actif, elle est bien plus
saturée que l'hémoglobine à la même pression, si bien que le dioxygène
passe de l'hémoglobine à la myoglobine.

\textbf{5.} $n = 1$.

\textbf{6.} La fixation sur un \omterm{def:b3:bioinorganic:haem}{hème} fait entrer son fer dans le plan et
tire l'histidine et son hélice~; les sous-unités basculent ensemble vers
une forme de plus haute affinité pour les autres sites.

\textbf{7.} $\qty{150}{g/L}/\qty{64500}{g/mol} = \qty{2.33e-3}{mol/L}$.

\textbf{8.} $4 \times \num{2.33e-3} = \qty{9.30e-3}{mol/L}$.

\textbf{9.} $\qty{9.30}{mmol/L} \times 0.972 = \qty{9.04}{mmol}$.

\textbf{10.} $\qty{9.30}{mmol/L} \times (0.972 - 0.724) = \qty{2.31}{mmol}$.

\textbf{11.} $\qty{2.31e-3}{mol} \times \qty{22.41}{L/mol} = \qty{52}{mL}$.

\textbf{12.} $0.248/0.972 = 26$\,\%.

\textbf{13.} $5.0^{2.7}/(4.0^{2.7} + 5.0^{2.7}) = 0.646$.

\textbf{14.} $\qty{9.30}{mmol/L} \times (0.972 - 0.646) = \qty{3.03}{mmol}$.

\textbf{15.} $3.03/2.31 = 1.31$~: trente pour cent de plus.

\textbf{16.} Le dioxyde de carbone produit par la respiration, hydraté en
acide carbonique par l'anhydrase carbonique, et l'acide lactique lors d'un
effort intense.

\textbf{17.} $\qty{3.03}{mmol/L} \times \qty{5.0}{L/min} = \qty{15.1}{mmol/min}$.

\textbf{18.} $\qty{15.1e-3}{mol/min} \times \qty{22.41}{L/mol} = \qty{0.34}{L/min}$.

\textbf{19.} $230 \times \num{100e-6}/0.2095 = 0.110$.

\textbf{20.} $0.110/1.110 = 0.099$~: environ 10\,\% des sites.

\textbf{21.} Les sites qui portent encore du dioxygène, dans une molécule
qui porte aussi du CO, sont sous la forme de haute affinité~: la courbe se
décale vers la gauche et ils libèrent moins de leur dioxygène dans les
tissus.

\textbf{22.} D'après la relation de compétition, le rapport $[\mathrm{HbCO}]/[\mathrm{HbO_2}]$ diminue proportionnellement à $p(\ce{O2})$~:
respirer du dioxygène pur, à une pression presque cinq fois plus élevée que
dans l'air, déplace plus vite le monoxyde de carbone.

\textbf{23.} En gênant la fixation linéaire de CO et en formant une liaison hydrogène avec $\ce{O2}$, elle maintient $M$ à
quelques centaines au lieu de la valeur bien plus grande de l'\omterm{def:b3:bioinorganic:haem}{hème} libre~;
sinon, même le monoxyde de carbone produit par l'organisme bloquerait une
grande partie de l'hémoglobine.

\textbf{24.} Au repos, le sang délivre environ \qty{0.34}{L} de dioxygène par minute.
\end{solution}
